% Options for packages loaded elsewhere
\PassOptionsToPackage{unicode}{hyperref}
\PassOptionsToPackage{hyphens}{url}
\PassOptionsToPackage{dvipsnames,svgnames,x11names}{xcolor}
\PassOptionsToPackage{space}{xeCJK}
\documentclass[
]{article}
\usepackage{xcolor}
\usepackage[margin=2.5cm]{geometry}
\usepackage{amsmath,amssymb}
\setcounter{secnumdepth}{-\maxdimen} % remove section numbering
\usepackage{iftex}
\ifPDFTeX
  \usepackage[T1]{fontenc}
  \usepackage[utf8]{inputenc}
  \usepackage{textcomp} % provide euro and other symbols
\else % if luatex or xetex
  \usepackage{unicode-math} % this also loads fontspec
  \defaultfontfeatures{Scale=MatchLowercase}
  \defaultfontfeatures[\rmfamily]{Ligatures=TeX,Scale=1}
\fi
\usepackage{lmodern}
\ifPDFTeX\else
  % xetex/luatex font selection
  \setmainfont[]{DejaVu Sans}
  \setmonofont[]{DejaVu Sans Mono}
  \ifXeTeX
    \usepackage{xeCJK}
    \setCJKmainfont[]{Noto Sans CJK SC}
  \fi
  \ifLuaTeX
    \usepackage[]{luatexja-fontspec}
    \setmainjfont[]{Noto Sans CJK SC}
  \fi
\fi
% Use upquote if available, for straight quotes in verbatim environments
\IfFileExists{upquote.sty}{\usepackage{upquote}}{}
\IfFileExists{microtype.sty}{% use microtype if available
  \usepackage[]{microtype}
  \UseMicrotypeSet[protrusion]{basicmath} % disable protrusion for tt fonts
}{}
\makeatletter
\@ifundefined{KOMAClassName}{% if non-KOMA class
  \IfFileExists{parskip.sty}{%
    \usepackage{parskip}
  }{% else
    \setlength{\parindent}{0pt}
    \setlength{\parskip}{6pt plus 2pt minus 1pt}}
}{% if KOMA class
  \KOMAoptions{parskip=half}}
\makeatother
\usepackage{color}
\usepackage{fancyvrb}
\newcommand{\VerbBar}{|}
\newcommand{\VERB}{\Verb[commandchars=\\\{\}]}
\DefineVerbatimEnvironment{Highlighting}{Verbatim}{commandchars=\\\{\}}
% Add ',fontsize=\small' for more characters per line
\newenvironment{Shaded}{}{}
\newcommand{\AlertTok}[1]{\textcolor[rgb]{1.00,0.00,0.00}{\textbf{#1}}}
\newcommand{\AnnotationTok}[1]{\textcolor[rgb]{0.38,0.63,0.69}{\textbf{\textit{#1}}}}
\newcommand{\AttributeTok}[1]{\textcolor[rgb]{0.49,0.56,0.16}{#1}}
\newcommand{\BaseNTok}[1]{\textcolor[rgb]{0.25,0.63,0.44}{#1}}
\newcommand{\BuiltInTok}[1]{\textcolor[rgb]{0.00,0.50,0.00}{#1}}
\newcommand{\CharTok}[1]{\textcolor[rgb]{0.25,0.44,0.63}{#1}}
\newcommand{\CommentTok}[1]{\textcolor[rgb]{0.38,0.63,0.69}{\textit{#1}}}
\newcommand{\CommentVarTok}[1]{\textcolor[rgb]{0.38,0.63,0.69}{\textbf{\textit{#1}}}}
\newcommand{\ConstantTok}[1]{\textcolor[rgb]{0.53,0.00,0.00}{#1}}
\newcommand{\ControlFlowTok}[1]{\textcolor[rgb]{0.00,0.44,0.13}{\textbf{#1}}}
\newcommand{\DataTypeTok}[1]{\textcolor[rgb]{0.56,0.13,0.00}{#1}}
\newcommand{\DecValTok}[1]{\textcolor[rgb]{0.25,0.63,0.44}{#1}}
\newcommand{\DocumentationTok}[1]{\textcolor[rgb]{0.73,0.13,0.13}{\textit{#1}}}
\newcommand{\ErrorTok}[1]{\textcolor[rgb]{1.00,0.00,0.00}{\textbf{#1}}}
\newcommand{\ExtensionTok}[1]{#1}
\newcommand{\FloatTok}[1]{\textcolor[rgb]{0.25,0.63,0.44}{#1}}
\newcommand{\FunctionTok}[1]{\textcolor[rgb]{0.02,0.16,0.49}{#1}}
\newcommand{\ImportTok}[1]{\textcolor[rgb]{0.00,0.50,0.00}{\textbf{#1}}}
\newcommand{\InformationTok}[1]{\textcolor[rgb]{0.38,0.63,0.69}{\textbf{\textit{#1}}}}
\newcommand{\KeywordTok}[1]{\textcolor[rgb]{0.00,0.44,0.13}{\textbf{#1}}}
\newcommand{\NormalTok}[1]{#1}
\newcommand{\OperatorTok}[1]{\textcolor[rgb]{0.40,0.40,0.40}{#1}}
\newcommand{\OtherTok}[1]{\textcolor[rgb]{0.00,0.44,0.13}{#1}}
\newcommand{\PreprocessorTok}[1]{\textcolor[rgb]{0.74,0.48,0.00}{#1}}
\newcommand{\RegionMarkerTok}[1]{#1}
\newcommand{\SpecialCharTok}[1]{\textcolor[rgb]{0.25,0.44,0.63}{#1}}
\newcommand{\SpecialStringTok}[1]{\textcolor[rgb]{0.73,0.40,0.53}{#1}}
\newcommand{\StringTok}[1]{\textcolor[rgb]{0.25,0.44,0.63}{#1}}
\newcommand{\VariableTok}[1]{\textcolor[rgb]{0.10,0.09,0.49}{#1}}
\newcommand{\VerbatimStringTok}[1]{\textcolor[rgb]{0.25,0.44,0.63}{#1}}
\newcommand{\WarningTok}[1]{\textcolor[rgb]{0.38,0.63,0.69}{\textbf{\textit{#1}}}}
\usepackage{longtable,booktabs,array}
\usepackage{caption}
\captionsetup[table]{skip=6pt}
\newcounter{none} % for unnumbered tables
\usepackage{calc} % for calculating minipage widths
% Correct order of tables after \paragraph or \subparagraph
\usepackage{etoolbox}
\makeatletter
\patchcmd\longtable{\par}{\if@noskipsec\mbox{}\fi\par}{}{}
\makeatother
% Allow footnotes in longtable head/foot
\IfFileExists{footnotehyper.sty}{\usepackage{footnotehyper}}{\usepackage{footnote}}
\makesavenoteenv{longtable}
\setlength{\emergencystretch}{3em} % prevent overfull lines
\providecommand{\tightlist}{%
  \setlength{\itemsep}{0pt}\setlength{\parskip}{0pt}}
\usepackage{bookmark}
\IfFileExists{xurl.sty}{\usepackage{xurl}}{} % add URL line breaks if available
\urlstyle{same}
% fallback for those not using the hyperref driver hyperxmp:
\makeatletter
\@ifundefined{xmpquote}{\newcommand{\xmpquote}[1]{#1}}{}
\makeatother
\hypersetup{
  colorlinks=true,
  linkcolor={Maroon},
  filecolor={Maroon},
  citecolor={Blue},
  urlcolor={Blue},
  pdfcreator={LaTeX via pandoc}}

\author{}
\date{}

\begin{document}

\section{筑基篇课后习题
I：内容详细对照（与筑基篇原文逐项比对）}\label{ux7b51ux57faux7bc7ux8bfeux540eux4e60ux9898-iux5185ux5bb9ux8be6ux7ec6ux5bf9ux7167ux4e0eux7b51ux57faux7bc7ux539fux6587ux9010ux9879ux6bd4ux5bf9}

\begin{quote}
\textbf{声明 (Declaration)}: 本文\textbf{不代表任何数学结论}, 仅为基于
Lean 4 形式化的一种\textbf{实验观测报告} (对照文档)。

{\def\LTcaptype{none} % do not increment counter
\begin{longtable}[]{@{}
  >{\raggedright\arraybackslash}p{(\linewidth - 6\tabcolsep) * \real{0.2500}}
  >{\raggedright\arraybackslash}p{(\linewidth - 6\tabcolsep) * \real{0.2500}}
  >{\raggedright\arraybackslash}p{(\linewidth - 6\tabcolsep) * \real{0.2500}}
  >{\raggedright\arraybackslash}p{(\linewidth - 6\tabcolsep) * \real{0.2500}}@{}}
\toprule\noalign{}
\begin{minipage}[b]{\linewidth}\raggedright
习题
\end{minipage} & \begin{minipage}[b]{\linewidth}\raggedright
主题
\end{minipage} & \begin{minipage}[b]{\linewidth}\raggedright
Lean 形式化
\end{minipage} & \begin{minipage}[b]{\linewidth}\raggedright
DOI
\end{minipage} \\
\midrule\noalign{}
\endhead
\bottomrule\noalign{}
\endlastfoot
exercise\_i\_crossref & 习题 I 内容详细对照 & PatPvsNPExercise.lean &
10.5281/zenodo.21917523 \\
\end{longtable}
}
\end{quote}

\textbf{Detailed Content Cross-Reference: Exercise I vs the
Foundation-Building Chapter}

\begin{quote}
本对照文档服务对象：\texttt{pat\_pnp\_exercise\_ZH.md}（R157，4
新定理）。逐项核对习题每个论点的锚定来源------题面（R152/R153 侧写）→
解答（R157 展开）→ 源定理（筑基篇 Lean
定理），三列对齐。全部源定理可复验：\texttt{formal/Formal/Toolkit/*.lean}。
\end{quote}

\emph{2026-08-13 · 对照文档 · 无新增 claim（仅比对，不产生新定理）}

\begin{quote}
外部文献对照（全字段书目 + 验证记录 + 健在/已故标注，85+
条）见\texttt{../shared/pat\_pnp\_shared\_references\_ZH.md}（共享引用清单）（引用清单），并已同步进
\texttt{literature/references.bib}。
\end{quote}

\begin{center}\rule{0.5\linewidth}{0.5pt}\end{center}

\subsection{一、对照总览}\label{ux4e00ux5bf9ux7167ux603bux89c8}

{\def\LTcaptype{none} % do not increment counter
\begin{longtable}[]{@{}
  >{\raggedright\arraybackslash}p{(\linewidth - 8\tabcolsep) * \real{0.2000}}
  >{\raggedright\arraybackslash}p{(\linewidth - 8\tabcolsep) * \real{0.2000}}
  >{\raggedright\arraybackslash}p{(\linewidth - 8\tabcolsep) * \real{0.2000}}
  >{\raggedright\arraybackslash}p{(\linewidth - 8\tabcolsep) * \real{0.2000}}
  >{\raggedright\arraybackslash}p{(\linewidth - 8\tabcolsep) * \real{0.2000}}@{}}
\toprule\noalign{}
\begin{minipage}[b]{\linewidth}\raggedright
习题章节
\end{minipage} & \begin{minipage}[b]{\linewidth}\raggedright
习题论点
\end{minipage} & \begin{minipage}[b]{\linewidth}\raggedright
锚定 claim
\end{minipage} & \begin{minipage}[b]{\linewidth}\raggedright
锚定 Lean 文件
\end{minipage} & \begin{minipage}[b]{\linewidth}\raggedright
对照结论
\end{minipage} \\
\midrule\noalign{}
\endhead
\bottomrule\noalign{}
\endlastfoot
零、总论 & P vs NP = 发散/收敛互锁对的投影 & R047 / R147 / R148 &
CausalityTime / InterlockIsomorphism & 论点 = R148
互锁同构的计算投影，无新增 \\
一、结构找回精确 & 验证 = 查表判等（双向） & R152 ② / RulerLookup /
RulerRevLookup & ConciseMagicTeaching & ★新定理 lookup\_iff\_value：R152
verification\_free 单向 → 双向 \\
二、模糊与结构等价 & witness ⟺ 表条目（双向） & R153 ④ / RulerLookup &
ConciseMagicTeaching / PatNondeterminism & ★新定理
witness\_iff\_table\_entry：R153 np\_witness\_in\_table 单向 → 双向 \\
三、★NP 模糊/结构对偶 & 模糊存在性 ∧ 结构判等 & R147 / R152 / R153 /
R150 & 本习题新定理 & ★新定理
np\_fuzzy\_structure\_duality：核心，前两定理组合 \\
四、全景 & N=外推 ∧ 验证=判等 ∧ P=锁定链 & R150 / R152 / R153 ① / R050 &
PatCountableInfinitPhaseUnification / PatNondeterminism & ★新定理
p\_np\_pat\_perspective：全景组合 \\
\end{longtable}
}

\textbf{一句话对照}：习题没有发明新结构------四个新定理都是筑基篇已有定理的\textbf{双向化}（单向→双向）与\textbf{组合化}（分离→对偶/全景）。这正是''用结构找回精确''在形式化层面的体现：锚定结构不变，论证收束。

\begin{center}\rule{0.5\linewidth}{0.5pt}\end{center}

\subsection{二、定理级对照（R157 新定理 ↔
源定理，逐证明步骤）}\label{ux4e8cux5b9aux7406ux7ea7ux5bf9ux7167r157-ux65b0ux5b9aux7406-ux6e90ux5b9aux7406ux9010ux8bc1ux660eux6b65ux9aa4}

\subsubsection{\texorpdfstring{2.1 \texttt{lookup\_iff\_value} --- 验证
=
查表判等（双向）}{2.1 lookup\_iff\_value --- 验证 = 查表判等（双向）}}\label{lookup_iff_value-ux9a8cux8bc1-ux67e5ux8868ux5224ux7b49ux53ccux5411}

\begin{Shaded}
\begin{Highlighting}[]
\NormalTok{theorem lookup\_iff\_value \{D E\} [Fintype D] [DecidableEq D] [DecidableEq E]}
\NormalTok{    (f : D → E) (x : D) (y : E) : y = f x ↔ (x, y) ∈ makeTable f}
\end{Highlighting}
\end{Shaded}

{\def\LTcaptype{none} % do not increment counter
\begin{longtable}[]{@{}
  >{\raggedright\arraybackslash}p{(\linewidth - 6\tabcolsep) * \real{0.2500}}
  >{\raggedright\arraybackslash}p{(\linewidth - 6\tabcolsep) * \real{0.2500}}
  >{\raggedright\arraybackslash}p{(\linewidth - 6\tabcolsep) * \real{0.2500}}
  >{\raggedright\arraybackslash}p{(\linewidth - 6\tabcolsep) * \real{0.2500}}@{}}
\toprule\noalign{}
\begin{minipage}[b]{\linewidth}\raggedright
证明步骤
\end{minipage} & \begin{minipage}[b]{\linewidth}\raggedright
使用的源定理
\end{minipage} & \begin{minipage}[b]{\linewidth}\raggedright
来源
\end{minipage} & \begin{minipage}[b]{\linewidth}\raggedright
与原题面（R152）的关系
\end{minipage} \\
\midrule\noalign{}
\endhead
\bottomrule\noalign{}
\endlastfoot
正向 \texttt{y\ =\ f\ x\ ⟹\ 表条目} & \texttt{rw\ {[}h{]}} +
\texttt{lookup\_exists} & ConciseMagicTeaching（RulerLookup 表全量） &
R152 \texttt{verification\_free} 只给
\texttt{(x,\ f\ x)\ ∈\ makeTable\ f}（特例） \\
反向 \texttt{表条目\ ⟹\ y\ =\ f\ x} & \texttt{lookup\_correct} &
ConciseMagicTeaching（RulerLookup 表值 = 函数值） & R152
未给反向；\textbf{习题新增双向} \\
值域多态 \texttt{{[}DecidableEq\ E{]}} & --- & 习题修正（验证函数 Bool
值域，非 D→D） & 初版错写 D→D，已修正 \\
\end{longtable}
}

\textbf{对照结论}：R152 verification\_free 是 lookup\_iff\_value 在 y =
f x 时的特例（代入即得）；lookup\_iff\_value 是 RulerLookup
两个单引理（lookup\_exists/lookup\_correct）的 iff 组合。

\subsubsection{\texorpdfstring{2.2 \texttt{witness\_iff\_table\_entry}
--- witness ⟺
表条目（双向）}{2.2 witness\_iff\_table\_entry --- witness ⟺ 表条目（双向）}}\label{witness_iff_table_entry-witness-ux8868ux6761ux76eeux53ccux5411}

\begin{Shaded}
\begin{Highlighting}[]
\NormalTok{theorem witness\_iff\_table\_entry \{D\} [Fintype D] [DecidableEq D]}
\NormalTok{    (V : D → D → Bool) (x : D) :}
\NormalTok{    (∃ w, V x w = true) ↔ (∃ w, (w, true) ∈ makeTable (fun w =\textgreater{} V x w))}
\end{Highlighting}
\end{Shaded}

{\def\LTcaptype{none} % do not increment counter
\begin{longtable}[]{@{}
  >{\raggedright\arraybackslash}p{(\linewidth - 6\tabcolsep) * \real{0.2500}}
  >{\raggedright\arraybackslash}p{(\linewidth - 6\tabcolsep) * \real{0.2500}}
  >{\raggedright\arraybackslash}p{(\linewidth - 6\tabcolsep) * \real{0.2500}}
  >{\raggedright\arraybackslash}p{(\linewidth - 6\tabcolsep) * \real{0.2500}}@{}}
\toprule\noalign{}
\begin{minipage}[b]{\linewidth}\raggedright
证明步骤
\end{minipage} & \begin{minipage}[b]{\linewidth}\raggedright
使用的源定理
\end{minipage} & \begin{minipage}[b]{\linewidth}\raggedright
来源
\end{minipage} & \begin{minipage}[b]{\linewidth}\raggedright
与原题面（R153）的关系
\end{minipage} \\
\midrule\noalign{}
\endhead
\bottomrule\noalign{}
\endlastfoot
正向 \texttt{witness\ ⟹\ 表条目} & \texttt{rw\ {[}h{]}} +
\texttt{lookup\_exists} & ConciseMagicTeaching & 与 R153
\texttt{np\_witness\_in\_table} \textbf{完全相同}（习题复用） \\
反向 \texttt{表条目\ ⟹\ witness} & \texttt{lookup\_correct}
反向（\texttt{.symm}） & ConciseMagicTeaching & R153
未给反向；\textbf{习题新增双向} \\
\end{longtable}
}

\textbf{对照结论}：R153 np\_witness\_in\_table 是
witness\_iff\_table\_entry 的正向合取；反向由 lookup\_correct（表值 =
函数值 ⟹ true = V x w ⟹ V x w =
true）封死。\textbf{模糊（witness）与结构（表条目）等价}的论证核心是
RulerLookup 的 lookup\_correct。

\subsubsection{\texorpdfstring{2.3
★\texttt{np\_fuzzy\_structure\_duality} --- NP
的模糊/结构对偶（核心新定理）}{2.3 ★np\_fuzzy\_structure\_duality --- NP 的模糊/结构对偶（核心新定理）}}\label{np_fuzzy_structure_duality-np-ux7684ux6a21ux7ccaux7ed3ux6784ux5bf9ux5076ux6838ux5fc3ux65b0ux5b9aux7406}

\begin{Shaded}
\begin{Highlighting}[]
\NormalTok{theorem np\_fuzzy\_structure\_duality \{D\} [Fintype D] [DecidableEq D]}
\NormalTok{    (V : D → D → Bool) (x : D) (hw : ∃ w, V x w = true) :}
\NormalTok{    (∃ w, (w, true) ∈ makeTable (fun w =\textgreater{} V x w)) ∧}
\NormalTok{    (∀ w, V x w = true ↔ (w, true) ∈ makeTable (fun w =\textgreater{} V x w))}
\end{Highlighting}
\end{Shaded}

{\def\LTcaptype{none} % do not increment counter
\begin{longtable}[]{@{}
  >{\raggedright\arraybackslash}p{(\linewidth - 4\tabcolsep) * \real{0.3333}}
  >{\raggedright\arraybackslash}p{(\linewidth - 4\tabcolsep) * \real{0.3333}}
  >{\raggedright\arraybackslash}p{(\linewidth - 4\tabcolsep) * \real{0.3333}}@{}}
\toprule\noalign{}
\begin{minipage}[b]{\linewidth}\raggedright
证明步骤
\end{minipage} & \begin{minipage}[b]{\linewidth}\raggedright
使用的源定理
\end{minipage} & \begin{minipage}[b]{\linewidth}\raggedright
来源
\end{minipage} \\
\midrule\noalign{}
\endhead
\bottomrule\noalign{}
\endlastfoot
前半 \texttt{结构存在性} & \texttt{witness\_iff\_table\_entry\ .mp} &
习题 2.2（其反向即 R153+lookup\_correct） \\
后半 \texttt{逐点判等} & \texttt{lookup\_iff\_value}（逐点实例化） &
习题 2.1（RulerLookup 双向） \\
\end{longtable}
}

\textbf{对照结论}：本题核心定理无新锚------是 2.1 与 2.2
的机械组合，但语义是新的：\textbf{``模糊存在性 ⟹ 结构存在性 ∧
逐点判等''把 R153 的''NP 存在性 =
验证表条目''从单向断言升级为对偶结构}。对应论证：求解（前向，发散）与验证（后向，收敛）=
R147 互锁对（CausalityTime.causality\_pair\_reduces: (f-e)+(e-f)=0）。

\subsubsection{\texorpdfstring{2.4 \texttt{p\_np\_pat\_perspective} ---
全景（N=外推 ∧ 验证=判等 ∧
P=锁定链）}{2.4 p\_np\_pat\_perspective --- 全景（N=外推 ∧ 验证=判等 ∧ P=锁定链）}}\label{p_np_pat_perspective-ux5168ux666fnux5916ux63a8-ux9a8cux8bc1ux5224ux7b49-pux9501ux5b9aux94fe}

\begin{Shaded}
\begin{Highlighting}[]
\NormalTok{theorem p\_np\_pat\_perspective \{D\} [Fintype D] [DecidableEq D]}
\NormalTok{    (V : D → D → Bool) (x : D) (θ : ℝ) (hθ₁ : 0 ≤ θ) (hθ₂ : θ ≤ 2 * Real.pi) :}
\NormalTok{    (∀ ε, 0 \textless{} ε → ∃ y ∈ patGrid, |θ {-} y| ≤ ε) ∧}
\NormalTok{    (∀ w, V x w = true ↔ (w, true) ∈ makeTable (fun w =\textgreater{} V x w)) ∧}
\NormalTok{    (∀ d, Function.Injective (fun x : ℝ =\textgreater{} x + d))}
\end{Highlighting}
\end{Shaded}

{\def\LTcaptype{none} % do not increment counter
\begin{longtable}[]{@{}
  >{\raggedright\arraybackslash}p{(\linewidth - 6\tabcolsep) * \real{0.2500}}
  >{\raggedright\arraybackslash}p{(\linewidth - 6\tabcolsep) * \real{0.2500}}
  >{\raggedright\arraybackslash}p{(\linewidth - 6\tabcolsep) * \real{0.2500}}
  >{\raggedright\arraybackslash}p{(\linewidth - 6\tabcolsep) * \real{0.2500}}@{}}
\toprule\noalign{}
\begin{minipage}[b]{\linewidth}\raggedright
合取项
\end{minipage} & \begin{minipage}[b]{\linewidth}\raggedright
使用的源定理
\end{minipage} & \begin{minipage}[b]{\linewidth}\raggedright
来源
\end{minipage} & \begin{minipage}[b]{\linewidth}\raggedright
对应题面
\end{minipage} \\
\midrule\noalign{}
\endhead
\bottomrule\noalign{}
\endlastfoot
① N = 相位锁定外推 & \texttt{pat\_phase\_unification}（王氏定理） &
PatCountableInfinitPhaseUnification（R150） & R153
③（nondeterminism\_is\_phase\_locking\_extrapolation 同款） \\
② 验证 = 查表判等 & \texttt{lookup\_iff\_value} 逐点 & 习题 2.1 & R152
②（verification\_free 双向化） \\
③ P = 锁定方向唯一链 & \texttt{deterministic\_locked\_chain\_unique} &
PatNondeterminism（R153 ① / R050 机制） & R153 ① 直接复用 \\
\end{longtable}
}

\textbf{对照结论}：全景定理 = R150（王氏）⊕ R153 ①（P）⊕ 习题
2.1（验证）------筑基篇三根支柱在 P vs NP 上的合取。R153
nondeterminism\_pat\_perspective 是 (①②) 子集；习题全景新增 ③
验证判等项并锚定 P。

\begin{center}\rule{0.5\linewidth}{0.5pt}\end{center}

\subsection{三、题面对照：R152/R153 说了什么，R157
补了什么}\label{ux4e09ux9898ux9762ux5bf9ux7167r152r153-ux8bf4ux4e86ux4ec0ux4e48r157-ux8865ux4e86ux4ec0ux4e48}

{\def\LTcaptype{none} % do not increment counter
\begin{longtable}[]{@{}
  >{\raggedright\arraybackslash}p{(\linewidth - 6\tabcolsep) * \real{0.2500}}
  >{\raggedright\arraybackslash}p{(\linewidth - 6\tabcolsep) * \real{0.2500}}
  >{\raggedright\arraybackslash}p{(\linewidth - 6\tabcolsep) * \real{0.2500}}
  >{\raggedright\arraybackslash}p{(\linewidth - 6\tabcolsep) * \real{0.2500}}@{}}
\toprule\noalign{}
\begin{minipage}[b]{\linewidth}\raggedright
原题面 claim
\end{minipage} & \begin{minipage}[b]{\linewidth}\raggedright
原定理（数量）
\end{minipage} & \begin{minipage}[b]{\linewidth}\raggedright
R157 对照
\end{minipage} & \begin{minipage}[b]{\linewidth}\raggedright
增量性质
\end{minipage} \\
\midrule\noalign{}
\endhead
\bottomrule\noalign{}
\endlastfoot
R152 有限域上 P=NP 平凡 & finite\_domain\_P\_eq\_NP &
习题四论证引用，无新定理 & 复用 \\
R152 验证 = 可逆查表后向免费 & verification\_free（单向） &
lookup\_iff\_value（双向） & \textbf{双向化} \\
R152 求解/验证 = 发散/收敛互锁对 & solve\_verify\_dual &
习题三论证引用（causality\_pair\_reduces） & 复用（论证层） \\
R153 ① P = 锁定方向唯一链 & deterministic\_locked\_chain\_unique &
习题四 ③ 合取项直接复用 & 复用 \\
R153 ② N = 未锁定多路径 & nondeterministic\_multiple\_paths &
习题四论证引用（未锁 = 模糊） & 复用（论证层） \\
R153 ③ N = 相位锁定外推 &
nondeterminism\_is\_phase\_locking\_extrapolation & 习题四 ①
合取项（同源 R150） & 复用 \\
R153 ④ NP 存在性 = 验证表条目 & np\_witness\_in\_table（单向） &
witness\_iff\_table\_entry（双向） & \textbf{双向化} \\
R153 ⑤ N 的 Pat 视角组合 & nondeterminism\_pat\_perspective &
np\_fuzzy\_structure\_duality + p\_np\_pat\_perspective &
\textbf{组合化+对偶化} \\
\end{longtable}
}

\textbf{增量清单（R157 独有的 4 个定理）}： 1.
\texttt{lookup\_iff\_value}：验证判等双向（R152 未给反向）。 2.
\texttt{witness\_iff\_table\_entry}：witness⟺表条目双向（R153
未给反向）。 3.
\texttt{np\_fuzzy\_structure\_duality}：模糊/结构对偶（新组合，核心）。
4. \texttt{p\_np\_pat\_perspective}：全景合取（R150⊕R153①⊕习题①）。

\begin{center}\rule{0.5\linewidth}{0.5pt}\end{center}

\subsection{四、论证链对照（习题章节 ↔
筑基篇章节）}\label{ux56dbux8bbaux8bc1ux94feux5bf9ux7167ux4e60ux9898ux7ae0ux8282-ux7b51ux57faux7bc7ux7ae0ux8282}

{\def\LTcaptype{none} % do not increment counter
\begin{longtable}[]{@{}
  >{\raggedright\arraybackslash}p{(\linewidth - 4\tabcolsep) * \real{0.3333}}
  >{\raggedright\arraybackslash}p{(\linewidth - 4\tabcolsep) * \real{0.3333}}
  >{\raggedright\arraybackslash}p{(\linewidth - 4\tabcolsep) * \real{0.3333}}@{}}
\toprule\noalign{}
\begin{minipage}[b]{\linewidth}\raggedright
习题章节
\end{minipage} & \begin{minipage}[b]{\linewidth}\raggedright
论证步骤
\end{minipage} & \begin{minipage}[b]{\linewidth}\raggedright
对应筑基篇章节/claim
\end{minipage} \\
\midrule\noalign{}
\endhead
\bottomrule\noalign{}
\endlastfoot
零、总论 & 发散/收敛互锁对投影 & 筑基篇·十二（R147
因果与时间，causality\_pair\_reduces）；筑基篇·十三（R148 互锁同构） \\
一、①表=结构 & 有限域函数=预计算表 &
筑基篇·RulerLookup（function\_is\_table）；R057 一切皆表 \\
一、②验证不重算 & 查表判等 & 筑基篇·R152 ②（verification\_free） \\
二、①单向已有 & R153 np\_witness\_in\_table & 筑基篇·十八（R153 ④） \\
二、②反向封死 & lookup\_correct 表值=函数值 &
筑基篇·RulerLookup（lookup\_correct） \\
三、①模糊换速度 & N 存在性由相位锁定外推 & 筑基篇·十七·十八（R150
王氏定理 / R153 ③） \\
三、②结构找回精确 & 表判等还原真值 &
筑基篇·RulerLookup/RulerRevLookup \\
三、③互锁对 & 发散端存在性/收敛端精确性 & 筑基篇·十二（R147）；R047
发散/收敛同一对称性 \\
四、①N=外推 & pat\_phase\_unification & 筑基篇·十六（R150
王氏相位锁定性定理） \\
四、②验证=判等 & lookup\_iff\_value & 筑基篇·十七（R152）+
RulerLookup \\
四、③P=锁定链 & deterministic\_locked\_chain\_unique & 筑基篇·十八（R153
①）；R050 机制；R136 方向声明 \\
四、④有限域平凡 & 表全量 ⟹ 模糊不必要 & 筑基篇·十七（R152 ①
finite\_domain\_P\_eq\_NP） \\
四、⑤诚实边界 & 非 P≠NP 判定 & 筑基篇·十七（R152 诚实边界原句延续） \\
\end{longtable}
}

\textbf{结论}：习题全部论证步骤都能在筑基篇找到章节级落点，无悬空引用；新增内容只发生在''双向化/组合化''层面。

\begin{center}\rule{0.5\linewidth}{0.5pt}\end{center}

\subsection{五、术语对照}\label{ux4e94ux672fux8bedux5bf9ux7167}

{\def\LTcaptype{none} % do not increment counter
\begin{longtable}[]{@{}
  >{\raggedright\arraybackslash}p{(\linewidth - 4\tabcolsep) * \real{0.3333}}
  >{\raggedright\arraybackslash}p{(\linewidth - 4\tabcolsep) * \real{0.3333}}
  >{\raggedright\arraybackslash}p{(\linewidth - 4\tabcolsep) * \real{0.3333}}@{}}
\toprule\noalign{}
\begin{minipage}[b]{\linewidth}\raggedright
习题术语
\end{minipage} & \begin{minipage}[b]{\linewidth}\raggedright
筑基篇对应概念
\end{minipage} & \begin{minipage}[b]{\linewidth}\raggedright
锚定
\end{minipage} \\
\midrule\noalign{}
\endhead
\bottomrule\noalign{}
\endlastfoot
模糊（N 的存在性） & 未锁定相位 & R140
single\_symmetry\_underdetermines（未锁定 ⟹ 位置不唯一）；R138 未锁定 =
自指坍缩 \\
速度（换来的） & 不穷举的存在性 & R150
王氏定理（任意精度统一，不依赖搜索路径） \\
结构（表） & 预计算表 / 一切皆表 & RulerLookup（makeTable）；R057
存储⟷计算同构 \\
精确（找回的） & 判等还原真值 & lookup\_correct（表值 = 函数值） \\
锁定链（P） & 方向锁定的唯一链 &
R050（迭代单射）；R136（成对一次性声明）；R137（pat n = pat0+n·d） \\
互锁对 & 求解/验证 = 发散/收敛 &
R147（causality\_pair\_reduces）；R047（发散/收敛同一对称性） \\
\end{longtable}
}

\begin{center}\rule{0.5\linewidth}{0.5pt}\end{center}

\subsection{六、诚实边界对照}\label{ux516dux8bdaux5b9eux8fb9ux754cux5bf9ux7167}

{\def\LTcaptype{none} % do not increment counter
\begin{longtable}[]{@{}
  >{\raggedright\arraybackslash}p{(\linewidth - 6\tabcolsep) * \real{0.2500}}
  >{\raggedright\arraybackslash}p{(\linewidth - 6\tabcolsep) * \real{0.2500}}
  >{\raggedright\arraybackslash}p{(\linewidth - 6\tabcolsep) * \real{0.2500}}
  >{\raggedright\arraybackslash}p{(\linewidth - 6\tabcolsep) * \real{0.2500}}@{}}
\toprule\noalign{}
\begin{minipage}[b]{\linewidth}\raggedright
边界
\end{minipage} & \begin{minipage}[b]{\linewidth}\raggedright
R152 原文
\end{minipage} & \begin{minipage}[b]{\linewidth}\raggedright
R157 延续
\end{minipage} & \begin{minipage}[b]{\linewidth}\raggedright
对照
\end{minipage} \\
\midrule\noalign{}
\endhead
\bottomrule\noalign{}
\endlastfoot
判定边界 & 完整 P≠NP 需要计算模型下界证明 & 习题仅交付结构侧写 &
一致，无突破 \\
定理边界 & 4 定理均为有限域/保留日志条件 & 4 新定理同样限 Fintype D &
一致（Fintype 前提继承） \\
引理边界 & R153 用户纠正：不用外部引理 & 习题全部锚本框架定理 & 一致，无
mathlib 独有引理（仅 ring/simp 基础） \\
\end{longtable}
}

\begin{center}\rule{0.5\linewidth}{0.5pt}\end{center}

\emph{筑基篇课后习题 I 对照文档 · 2026-08-13 · 无新增
claim，仅比对（源定理全部可复验于 formal/Formal/Toolkit/）}

\end{document}
